% !TEX encoding = UTF-8
% !TEX TS-program = pdflatex
% !TEX root = ../tesi.tex

%**************************************************************
\chapter{Processi e metodologie}
\label{cap:processi-metodologie}
%**************************************************************

\intro{In questo capitolo verranno descritte le principali tecnologie che ho approfondito durante l'esperienza di stage.}\\

%**************************************************************
\section{Tecnologie utilizzate}
Durante la mia esperienza di stage ho potuto prendere contatto con molte tecnologie differenti e che non avevo incontrato durante il corso di studi. Esse spaziano dall'ambiente server a quello client. L'approfondimento di tali tecnologie, specialmente nelle prime settimane di stage, è stato fondamentale per capire il funzionamento del software da modificare \\
\subsection{Java}
\begin{figure}[h]
	\centering
	\includegraphics[height = 4 cm]{java-logo}
	\caption{Logo di Java}
\end{figure}
Java è uno dei linguaggi di programmazione general-purpose più popolari al mondo. \\
Rilasciato nel 1995, riprende molta della sintassi dal linguaggio C e C++, ma spostando parte delle responsabilità prima lasciate ai programmatori, come la gestione della memoria, alla \emph{Java Virtual Machine (JVM)}: una macchina virtuale su chi viene eseguito tutto il codice Java.\\
La presenza di una macchina virtuale crea un livello aggiuntivo tra il sistema operativo e l' \emph{IDE} di sviluppo: Questo permette al linguaggio Java di aderire al principio \emph{Write Once Run Anywhere (WORA)}: Il codice compilato (bytecode) può essere eseguito da qualsisasi computer provvisto di una JVM, senza necessitare di ricompilazioni e perfino adattamenti del codice sorgente in base al sistema operativo presente sul computer.
Java è inoltre molto utilizzato per eseguire applicazioni client-server, cioè applicazioni (Client) che per ottentere i dati necessari si appoggiano ad un fornitore (Server), che recupera per loro i dati necessari, in maniera del tutto trasparente al client e quindi all'utente.\\
Java dispone di una libreria proprietaria specializzata, ma in JGalileoCRM tale libreria è sostituita da \emph{Liferay}.\\ %TODO occhio che java.net è chiusa
\subsection{Liferay}
\begin{figure}[h]
	\centering
	\includegraphics[height = 4 cm]{Logo-Liferay}
	\caption{Logo di Liferay}
\end{figure}
Liferay è una tecnologia \emph{enterprise portal} open source, realizzato in Java.\\
Per enterprise portal si intende un sistema informatico evoluto, in grado di integrare informazioni, processi e persone allo scopo di fornire valore aggiunto in termini di:
\begin{itemize}
	\item Gestione del \emph{Single Sign On};
	\item Semplice personalizzazioni ad hoc per ogni cliente;
	\item Analisi delle pagine, in termini di accessi, click, download molto semplice;
	\item Integrazione tra funzionalità e dati di diversi sistemi in nuove componenti definite portlets. 
\end{itemize}
Le portlets sono il cuore del sistema di Liferay. Esse infatti permettono di concentrare lo sviluppo solamente sulla gestione della funzionalità principale, lasciando a liferay la gestione degli accessi, dei menù di navigazione e degli altri componenti globali dell'applicazione.\\

\subsection{GWT}
\begin{figure}[h]
	\centering
	\includegraphics[height = 4 cm]{gwt-logo}
	\caption{Logo di Google Web Toolkit}
\end{figure} 
GWT, acronimo di \emph{Google Web Toolkit} è un insieme di tool open source che permette la creazione ed il mantenimento di complesse applicazioni front-end JavaScript scritte in Java.\\
GWT infatti si occupa di tradurre il codice scritto in Java in codice JavaScript, interpretabile da tutti i moderni browsers Tutto il codice Java può essere compilato ed eseguito grazie ai file Ant inclusi.
Ant è un progetto open source della Apache foundation, volto ad automatizzare il processo di build. è simile al comando make di Unix, ma scritto in Java. \\ %TODO make come codice
Il plug-in di Google per Eclipse (IDE in uso presso Sanmarco Informatica) è molto comleto, offrendo la possibilità di creare progetti, farne il debug, il testing, meccanismi di validazione e di controllo della sintassi.\\
\subsection{Apache Tomcat}
\begin{figure}[h]
	\centering
	\includegraphics[height = 4 cm]{apache_tomcat-card}
	\caption{Logo di Apache Tomcat}
\end{figure}
Apache Tomcat è un web server, cioè un' applicazione web che, in esecuzione su di un server, gestisce le richieste di trasferimento tra le varie pagine web di un client.\\
In particolare, Tomcat opera utilizzando servlet, cioè oggetti scritti in Java molto usati nella generazione di pagine web dinamiche.
Queste servlet, in JGalileoCRM non vengono direttamente scritte in Java, ma sono scritte con \emph{JavaServer Pages (JSP)}:\\
JSP è una tecnologia di programmazione web, scritta in Java, che si basa sull'uso di speciali tag all'interno di una pagina HTML, con cui possono essere chiamate funzioni scritte in linguaggio Java o JavaScript.
A runtime, le pagine JSP vengono tradotte automaticamente in servlet utilizzabili da Tomcat. \\
Il vantaggio principale di questa tecnologia rispetto, ad esempio, a \gls{PHP}, consiste nel poter scrivere tutto il codice, frontend e backend in un solo linguaggio di programmazione: Java. \\
Inoltre permette la creazioni di applicazioni web dinamiche, rispetto all'utilizzo di codice HTML statico.
\subsection{Databases}
Un database, o base di dati, è un insieme di dati omogeneo memorizzato in un elaboratore, che può essere interrogato attraverso uno dei linguaggi di interrogazione esistenti, con lo scopo di ottenere una parte dei dati memorizzati.
JGalileoCRM utilizza due databases per gestire i propri dati:
\subsubsection{Database SQL}
Il primo database utilizzato dal software sopracitato è un comune database \gls{SQL}: esso è un database relazionale, cioè opera creando tabelle che vengono messe in relazione tra loro tramite gli attributi inseriti nellle tabelle stesse.\\
Questo database è utilizzato per la gestione di tutti i dati diversi dai dati personali dei clienti e dei contatti. %TODO spiegare meglio
\subsubsection{Database AS400}
AS/400 (Application System 400)è un è un minicomputer sviluppato a partire dal 1988 dall’IBM per usi prevalentemente aziendali, come supporto del sistema informativo gestionale.\\
Il punto di forza sta nel database integrato con il sistema operativo, che permette la gestione dei dati attraverso una suite di programmi e librerie altamente specializzati. Attualmente consente di eseguire tutte le operazioni disponibili su di un comune server web, grazie ai continui aggiornamenti che hanno portato, ad esempio, all'installazione di PHP direttamente a livello di sistema operativo.
Questo database viene usato, in Sanmarco Informatica, per contenere tutti i dati degli utenti, sia interni che esterni, nonchè i dati relativi ai contatti %TODO meglio, chiedere a luca il tipo di DB
\subsection{Hibernate}
Hibernate è una piattaforma \gls{middleware} open-source per gestire la persistenza dei dati in un database.\\
è largamente utilizzato per la gestione dei dati di applicazioni web scritte in Java, in quanto i dati vengono rappresentati attraverso degli oggetti Java chiamati \emph{entità}.
Un'entità hibernate si sostituisce logicamente alla tabella reale del database. %todo meglio
\subsection{JavaScript}
\subsubsection{Analisi delle alternative}
\subsubsection{React}
\subsubsection{Angular}
\subsubsection{Vue.JS}
\subsubsection{Scelta finale}